\documentclass[12pt, twoside]{article}
%\documentclass[12pt, twoside]{article}
\usepackage[letterpaper, margin=1in, headsep=0.2in]{geometry}
\setlength{\headheight}{0.6in}
%\usepackage[english]{babel}
\usepackage[utf8]{inputenc}
\usepackage{microtype}
\usepackage{amsmath}
\usepackage{amssymb}
%\usepackage{amsfonts}
\usepackage[nomessages]{fp} %\FPeval{\var-name}{2*sin(pi/6)}
\usepackage{siunitx} %units in math. eg 20\milli\meter
\usepackage{yhmath} % for arcs, overparenth command
\usepackage{tikz} %graphics
\usetikzlibrary{quotes, angles, arrows, arrows.meta}
\usepackage{pgfplots}
\usepgfplotslibrary{fillbetween}
\pgfplotsset{compat=1.16}
\usepackage{graphicx} %consider setting \graphicspath{{images/}}
\usepackage{parskip} %no paragraph indent
\usepackage{enumitem}
\usepackage{multicol}
\usepackage{venndiagram}

\usepackage{fancyhdr}
\pagestyle{fancy}
\fancyhf{}
\renewcommand{\headrulewidth}{0pt} % disable the underline of the header
\raggedbottom
\hfuzz=2mm %suppresses overfull box warnings

\usepackage{hyperref}

\fancyhead[LE]{\thepage}
\fancyhead[RO]{Markscheme}
\fancyhead[LO]{La Scuola d'Italia / Dr.\ Huson / IB Math 12 \\* 14 October 2026}

\begin{document}
\subsubsection*{1.11 Quiz: Statistics --- Markscheme}

\begin{enumerate}[itemsep=0.3cm]

%--- Problem 1: Grouped frequency table; continuous data; modal class; estimated mean and sd; stratified sample [7 marks] ---
\item[1.] [Maximum mark: 7]

Mid-interval values $5, 15, 25, 35, 45$ with frequencies $4, 12, 18, 11, 5$;
$n = 50$.

\emph{\textbf{GDC:} enter the mid-interval values in one list and the
frequencies in a second list, then run \texttt{One-Variable Statistics} with
the frequency list.}

\textbf{(a)} Continuous \hfill \textbf{A1}

\emph{Time is measured, not counted. Accept a reason such as ``time can
take any value in an interval''; no reason is required.}

\textbf{(b)} $20 \le t < 30$ \hfill \textbf{A1}

\emph{The class, not its frequency: ``$18$'' alone earns A0. Accept
``$20$ to $30$ minutes''. The mid-interval value $25$ alone earns A0.}

\textbf{(c)} $\bar{x} = \dfrac{4(5) + 12(15) + 18(25) + 11(35) + 5(45)}{50}$
\hfill \textbf{(M1)}

$= \dfrac{1260}{50} = 25.2$ minutes \hfill \textbf{A1}

\emph{M1 for using the mid-interval values with the frequencies, shown as a
sum or as the GDC lists. Using the lower class boundaries
($0, 10, 20, 30, 40$, giving $20.2$) or upper boundaries (giving $30.2$)
earns M0A0. The mean of the frequencies, $10$, earns no marks.}

\textbf{(d)} $\sigma = 10.9 \quad (10.8609\ldots)$ minutes \hfill \textbf{A1}

\emph{A1 for $10.9$ (3 s.f.) or the display with an ellipsis
($10.8609\ldots$); $10.86$ with no ellipsis earns A0 (markscheme-spec §21). Do not accept $s_{x} = 10.971\ldots$
($11.0$), the sample standard deviation.}

\textbf{(e)} $\dfrac{192}{600} \times 50$ \hfill \textbf{M1}

$= 16$ patients \hfill \textbf{A1}

\emph{M1 for the over-$65$ proportion of the registered patients applied to
$50$, however written (e.g.\ $\frac{50}{600} = \frac{1}{12}$, then
$\frac{192}{12}$).}

\newpage

%--- Problem 2: Five-number summary from raw data; IQR; low-outlier test; draw the box plot [9 marks] ---
\item[2.] [Maximum mark: 9]

Ordered: $8, 22, 24, 25, 26, 27, 28, \mathbf{28}, 29, 30, 31, 32, 33, 35,
36$ \quad ($n = 15$).

\emph{\textbf{GDC:} the TI-Nspire excludes the median from each half, so
$Q_{1}$ is the median of the lowest seven values and $Q_{3}$ the median of
the highest seven.}

\textbf{(a)} Median $= 28$ \hfill \textbf{A1}

$Q_{1} = 25$, \quad $Q_{3} = 32$ \hfill \textbf{A1}

\emph{Both quartiles are needed for the second A1. Accept $Q_{1} = 25.5$ and
$Q_{3} = 31.5$ from a hand method that includes the median in each half;
feedback: the GDC convention is the one to use.}

\textbf{(b)} $\text{IQR} = 32 - 25 = 7$ \hfill \textbf{A1}

\emph{FT: award A1 for $Q_{3} - Q_{1}$ from the candidate's own quartiles.
The range, $36 - 8 = 28$, earns no mark.}

\textbf{(c)} $Q_{1} - 1.5 \times \text{IQR} = 25 - 1.5 \times 7$
\hfill \textbf{M1}

$= 14.5$ \hfill \textbf{A1}

$8 < 14.5$, so $8$ is an outlier. \hfill \textbf{R1}

\emph{M1 for $1.5 \times$ their IQR subtracted from their $Q_{1}$. R1 for
the comparison with $8$ stated, not just the fence value. Subtracting
$1.5 \times \text{IQR}$ from the median or the mean earns no marks. The
upper fence, $32 + 10.5 = 42.5$, is not required. FT: award all three marks
from the candidate's own $Q_{1}$ and IQR, provided their fence is above $8$.}

\textbf{(d)}

\begin{center}
\begin{tikzpicture}
  \begin{scope}[x=0.35cm]
    \draw (5,0)--(41,0);
    \foreach \t in {5,6,...,41} \draw (\t,-0.08)--(\t,0.08);
    \foreach \t in {5,10,...,40} \node[below] at (\t,-0.12) {\footnotesize \t};
    \draw[thick] (25,0.6) rectangle (32,1.4);
    \draw[thick] (28,0.6)--(28,1.4);
    \draw[thick] (22,1.0)--(25,1.0);
    \draw[thick] (32,1.0)--(36,1.0);
    \draw[thick] (22,0.8)--(22,1.2);
    \draw[thick] (36,0.8)--(36,1.2);
    \node at (8,1.0) {$\times$};
  \end{scope}
\end{tikzpicture}
\end{center}

Box from $Q_{1} = 25$ to $Q_{3} = 32$ with the median line at $28$
\hfill \textbf{A1}

Whiskers to $22$ and $36$ \hfill \textbf{A1}

Outlier $8$ marked with a cross \hfill \textbf{A1}

\emph{Award each mark independently; plotted values within half a small
division ($\pm 0.5$) are accepted. A whisker drawn to $8$ or to $14.5$ loses
the second A1. FT: award the first A1 for a box drawn at the candidate's own
quartiles from (a).}

\newpage

%--- Problem 3: Read median, IQR, a count and a percentile from a CF curve; frequencies from the curve [9 marks] ---
\item[3.] [Maximum mark: 9]

$200$ fish; the curve passes through $(20, 30)$, $(24, 50)$, $(28, 75)$,
$(32, 100)$, $(40, 150)$ and $(44, 170)$.

\emph{Readings within one small square of the values below are accepted
(\,$\pm 2$ cm on the length axis, $\pm 5$ fish on the cumulative frequency
axis), and any value following from the candidate's own readings earns the
A1. Feedback: draw the reading lines on the graph.}

\textbf{(a)} Median: cumulative frequency $100$ gives $L = 32$ cm
\hfill \textbf{A1}

\textbf{(b)} $Q_{1} = 24$ (at $50$), \quad $Q_{3} = 40$ (at $150$)
\hfill \textbf{(M1)}

$\text{IQR} = 40 - 24 = 16$ cm \hfill \textbf{A1}

\emph{M1 for reading the lengths at cumulative frequencies $50$ and $150$
(or for either quartile correct). Reading at $25$ and $75$ (percent of
$100$) earns M0.}

\textbf{(c)} $200 - 170$ \hfill \textbf{(M1)}

$= 30$ fish \hfill \textbf{A1}

\emph{M1 for reading the cumulative frequency at $L = 44$ and subtracting
from $200$. The reading $170$ alone, as a final answer, earns M1A0.}

\textbf{(d)} $15\% \text{ of } 200 = 30$ \hfill \textbf{(M1)}

Cumulative frequency $30$ gives $L = 20$ cm \hfill \textbf{A1}

\emph{M1 for $30$ (or for reading the length at cumulative frequency $30$).
Reading at cumulative frequency $170$ (the longest $15\%$, $L = 44$) earns
M1A0. Reading at cumulative frequency $15$ (treating the total as $100$,
$L = 16$) earns M0A0.}

\textbf{(e)} $p = 75 - 30 = 45$ \hfill \textbf{A1}

$q = 170 - 75 = 95$ \hfill \textbf{A1}

\emph{The two marks are independent. Accept $q = 200 - 30 - 45 - 30 = 95$.
FT: award the A1 for $q$ from the candidate's own $p$ when $q$ is found by
subtraction from $200$.}

\newpage

%--- Problem 4: Pearson's r, characterizing correlation, y-on-x regression line, interpreting the gradient, interpolation [8 marks] ---
\item[4.] [Maximum mark: 8]

$n = 8$, \quad $\bar{x} = 5.5$, \quad $\bar{y} = 18.875$.

\emph{\textbf{GDC:} enter $x$ and $y$ in two lists and run
\texttt{Linear Regression (mx+b)}.}

\textbf{(a)} $r = -0.973 \quad (-0.97297\ldots)$ \hfill \textbf{A1}

\emph{A1 for $-0.973$ (3 s.f.) or the display with an ellipsis
($-0.97297\ldots$); $-0.9730$ or a truncated $-0.972$ earns A0. The negative sign is required. Do not accept $r^{2} = 0.947$, which
sits beside $r$ on the regression screen.}

\textbf{(b)} Strong, negative (linear) correlation \hfill \textbf{A1A1}

\emph{A1 for ``negative'', A1 for ``strong'' (accept ``very strong''); the
two marks are independent. ``Moderate'' or ``weak'' earns A0 for the
strength. FT from the candidate's $r$ in (a).}

\textbf{(c)} $y = -1.84x + 29.0$ \hfill \textbf{A1A1}

\emph{A1 for $a = -1.84$, A1 for $b = 29.0$; the two marks are independent.
Each coefficient to 3 s.f., or as the display with an ellipsis ($a =
-1.8397\ldots$, $b = 28.993\ldots$); $a = -1.840$ or $b = 28.99$ with no
ellipsis earns A0 for that coefficient (markscheme-spec §21). $b = 29$ is 2 s.f.: A0 for $b$. The reversed line $x = -0.515y + 15.2$ earns A0A0.}

\textbf{(d)} For each additional year of age, the price decreases by about
$1.84$ thousand dollars (about \$1840). \hfill \textbf{A1}

\emph{The answer must give the rate (about $1.8$ thousand dollars) per year
and the direction (decrease), in context. ``Older cars are cheaper''
restates the correlation without the rate and earns A0. FT from the
candidate's $a$ in (c).}

\textbf{(e)} $y = -1.84(7) + 29.0$ \hfill \textbf{(M1)}

$= 16.1 \quad (16.115\ldots)$ thousand dollars \hfill \textbf{A1}

\emph{A1 for $16.1$ (3 s.f.), \$16\,100, or the display with an ellipsis
($16.115\ldots$), from any route: full-precision coefficients or the
rounded line ($-1.84(7) + 29.0 = 16.12$) both round to $16.1$. A final
answer of $16.12$ with no ellipsis earns A0 (markscheme-spec §21). FT from the
candidate's line in (c). \textbf{Feedback:} store the coefficients in the GDC
(or use the regression function \texttt{f1(7)}) and round only the final
answer.}

\end{enumerate}
\end{document}
